% Part: incompleteness
% Chapter: representability-in-q
% Section: c-representable

\documentclass[../../../include/open-logic-section]{subfiles}

\begin{document}

\olfileid{inc}{req}{cre}
\olsection{$C$లోని ప్రమేయాలను $\Th{Q}$లో ప్రతినిధీకరించవచ్చు}

$C$లోని ప్రతి ప్రమేయాన్ని $\Th{Q}$లో ప్రతినిధీకరించవచ్చని చూపాలి. చివరకు, $C$లోని ప్రతి $k$-స్థానిక ప్రమేయం $f(x_0,\dots,x_{k-1})$కు దాన్ని ప్రతినిధీకరించే $!A_f(x_0,\dots,x_{k-1},y)$ సూత్రాన్ని ఎలా కేటాయించాలో చూపాలి.

\begin{lem}
సహజ సంఖ్యలు $n$, $m$కు $n \neq m$ అయితే, $Q \Proves \num
n \neq \num m$.
\end{lem}

\begin{proof}
ప్రతి $m$కు $n \neq m$ అయితే $Q \Proves \eq/[\num n][\num m]$ అని $n$పై ఆగమన పద్ధతిలో చూపుదాం.

ఆధార సందర్భంలో $n = 0$. $m$, $0$కు సమానం కాకపోతే, ఏదో సహజ సంఖ్య $k$కు $m = k +
1$. $\lforall[x][\eq/[0][x']]$ అని చెప్పే స్వీకృతం మనకు ఉంది. పరిమాణీకరణిక స్వీకృతం ద్వారా $x$ స్థానంలో $\num k$ను పెట్టి $\eq/[0][\num k']$ అని తేల్చవచ్చు. కానీ $\num k'$ అనేది $\num m$యే.

ఆగమన దశలో, ప్రతిపాదన $n$కు నిజమని అనుకొని $n+1$ను పరిశీలిద్దాం. $m$ ఏ సహజ సంఖ్యైనా అనుకుందాం. రెండు అవకాశాలున్నాయి: $m = 0$ లేదా ఏదో $k$కు $m = k+1$. మొదటి సందర్భాన్ని పైనిలాగే నిర్వహించవచ్చు. రెండో సందర్భంలో $n+1 \neq
k+1$ అనుకుందాం. అప్పుడు $n \neq k$. $n$కు ఆగమన పరికల్పన ప్రకారం $\Th{Q} \Proves \eq/[\num n][\num k]$. $\lforall[x][\lforall[y][\eq[x'][y'] \lif \eq[x][y]]]$ అనే స్వీకృతం ఉంది. పరిమాణీకరణిక స్వీకృతాన్ని ఉపయోగిస్తే $\eq[\num n'][\num k'] \lif \eq[\num n][\num
  k]$ వస్తుంది. ప్రతిపాదనాత్మక తర్కం ద్వారా $\Th{Q}$లో $\eq/[\num n][\num k] \lif \eq/[\num n'][\num k']$ అని తేల్చవచ్చు. మోడస్ పోనెన్స్ ద్వారా $\eq/[\num n'][\num k']$ వస్తుంది; $\num k'$ అనేది $\num m$ కాబట్టి ఇదే కావలసినది.
\end{proof}

ఈ లెమ్మా చెప్పేది చాలా పరిమితమే: సారాంశంగా, వేర్వేరు సంఖ్యా సంకేతాలు వేర్వేరు వస్తువులను సూచిస్తాయని $\Th{Q}$ నిరూపించగలదని చెబుతుంది. ఉదాహరణకు $\Th{Q}$, $0'' \neq 0'''$ అని నిరూపిస్తుంది. అయితే ఇది సాధారణంగా నిజమని చూపడానికి కొంత జాగ్రత్త అవసరం. మనం ఆగమన పద్ధతిని వాడుతున్నా, అది $\Th{Q}$ \emph{బయట} చేసే ఆగమనమని కూడా గుర్తుంచుకోండి.

సున్నా, ఉత్తరవర్తి, కూడిక, గుణకారం, సమానత్వానికి లక్షణ ప్రమేయం, ప్రక్షేపణలను ప్రతినిధీకరించగలం. ప్రతి సందర్భంలో తగిన ప్రతినిధీకరణ సూత్రం చాలా సరళమే. ఉదాహరణకు సున్నాను ఈ సూత్రం ప్రతినిధీకరిస్తుంది:
\[
y = 0,
\]
ఉత్తరవర్తిని ఈ సూత్రం ప్రతినిధీకరిస్తుంది:
\[
x_0' = y,
\]
కూడికను ఈ సూత్రం ప్రతినిధీకరిస్తుంది:
\[
x_0 + x_1 = y.
\]
అసలు పని సంబంధిత వాక్యాలను $\Th{Q}$ నిరూపించగలదని చూపడంలో ఉంది. ఉదాహరణకు పైనున్న సూత్రం కూడికను ప్రతినిధీకరిస్తుందని చెప్పాలంటే, ప్రతి సహజ సంఖ్యల జత $m$, $n$కు $\Th{Q}$ కిందివి నిరూపిస్తుందని చూపాలి:
\[
\eq[\num n + \num m][\num {n+m}]
\]
మరియు
\[
\lforall[y][(\eq[(\num n + \num m)][y] \lif \eq[y][\num {n+m}])].
\]

సంయోజనం సంగతి ఏమిటి? $h$ను ఈ విధంగా నిర్వచించామనుకుందాం:
\[
h(x_0,\dots,x_{l-1}) = f(g_0(x_0,\dots,x_{l-1}), \dots,
g_{k-1}(x_0,\dots,x_{l-1})).
\]
ఇక్కడ $f,g_0,\dots,g_{k-1}$ ప్రమేయాలను వరుసగా ప్రతినిధీకరించే $!A_f, !A_{g_0}, \dots,
!A_{g_{k-1}}$ సూత్రాలను ముందే కనుగొన్నాం. అప్పుడు $h$ను ప్రతినిధీకరించే $!A_h$ సూత్రాన్ని, $!A_h(x_0,\dots,x_{l-1},y)$ను కిందివిధంగా నిర్వచించడం ద్వారా పొందవచ్చు:
\begin{multline*}
  \lexists[z_0\dots][\lexists[z_{k-1}][(!A_{g_0}(x_0,\dots,x_{l-1},z_0) \land
  \dots \land !A_{g_{k-1}}(x_0,\dots,x_{l-1},z_{k-1}) \land]]\\
  !A_f(z_0,\dots,z_{k-1},y)).
\end{multline*}

చివరగా అపరిమిత శోధనను పరిశీలిద్దాం. $g(x,\vec z)$ క్రమబద్ధమై, $\Th{Q}$లో ప్రతినిధీకరించబడుతుందని, దాని సూత్రం $!A_g(x,\vec
z,y)$ అని అనుకుందాం. $f$ను $f(\vec z) = \mu x \; g(x,\vec z)$గా నిర్వచిద్దాం. $f$ను ప్రతినిధీకరించే $!A_f(\vec z,y)$ సూత్రాన్ని కనుగొనాలనుకుంటున్నాం. సహజమైన ఎంపిక ఇది:
\[
!A_f(\vec z,y) \ident !A_g(y,\vec z,0) \land \lforall[w][(w < y
\lif \lnot !A_g(w,\vec z,0))].
\]

కొన్ని అదనపు లెమ్మాలను ఉపయోగించి ఇది పనిచేస్తుందని చూపవచ్చు; ఉదాహరణకు:

\begin{lem}
  ప్రతి చరరాశి $x$కు, ప్రతి సహజ సంఖ్య $n$కు, $\Th{Q}$ ఈ $\eq[(x'
  + \num n)][(x + \num n)']$ను నిరూపిస్తుంది.
\end{lem}

ఇది $\Th{Q}$, $\lforall[x][\lforall[y][(x' + y) = (x +y)']]$ను నిరూపిస్తుందని చెప్పడం కంటే బలహీనమైనదని మళ్లీ గుర్తుంచుకోవాలి (ఆ బలమైన వాదన అసత్యం).

\begin{proof}
ఎప్పటిలాగే నిరూపణ $n$పై ఆగమన పద్ధతిలో సాగుతుంది. ఆధార సందర్భంలో $n =
0$; $\Th{Q}$, $\eq[(x'+0)][(x+0)']$ను నిరూపిస్తుందని చూపాలి. మనకు ఇవి ఉన్నాయి:
\begin{eqnarray*}
& & \eq[(x' + 0)][x'] \quad \text{4వ స్వీకృతం నుంచి} \\
& & \eq[(x + 0)][x] \quad \text{4వ స్వీకృతం నుంచి} \\
& & \eq[(x+0)'][x'] \quad \text{2వ పంక్తి వల్ల} \\
& & \eq[(x' + 0)][(x + 0)'] \quad \text{1, 3 పంక్తుల వల్ల}
\end{eqnarray*}
ఆగమన దశలో $\Th{Q}$లో $\eq[(x' +
  \num n)][(x + \num n)']$ను వ్యుత్పాదించామని అనుకుందాం. $\num{n+1}$ అనేది $\num
n'$ కాబట్టి $\Th{Q}$, $\eq[(x' + \num n')][(x + \num
n')']$ను నిరూపిస్తుందని చూపాలి. ఈ సమానత్వాల గొలుసును $\Th{Q}$లో వ్యుత్పాదించవచ్చు:
\begin{eqnarray*}
(x' + \num n') & \eq & (x' + \num n)' \quad \text{5వ స్వీకృతం} \\
& \eq & (x + \num n')' \quad \text{ఆగమన పరికల్పన నుంచి}
\end{eqnarray*}
\end{proof}

\begin{lem}
\begin{enumerate}
\item $\Th{Q}$, $\lnot (x < \num 0)$ను నిరూపిస్తుంది.
\item ప్రతి సహజ సంఖ్య $n$కు $\Th{Q}$ కిందిదాన్ని నిరూపిస్తుంది:
\[
x < \num {n+1} \lif (\eq[x][\num 0] \lor \dots \lor \eq[x][\num n]).
\]
\end{enumerate}
\end{lem}

\begin{proof}
1ని, 2లో ఒక భాగాన్ని అనౌపచారికంగా చేద్దాం (అంటే ఔపచారిక \tetoken{వ్యుత్పత్తి}{!!{derivation}} ఎలా నిర్మించాలో సూచనలు మాత్రమే ఇస్తాం).

1వ భాగానికి, $<$ నిర్వచనం ప్రకారం $\Th{Q}$లో $\lnot
\lexists[y][\eq[(y'+x)][0]]$ను నిరూపించాలి. ఇది (మొదటిస్థాయి తర్కశాస్త్రపు స్వీకృతాలు, నియమాలు ఉపయోగించి) $\lforall[y][\eq/[(y' +
    x)][0]]$కు తుల్యం. ఆలోచన ఇది: $\eq[(y'+x)][0]$ అనుకుందాం. $x$, 0 అయితే $\eq[(y' + 0)][0]$ వస్తుంది. కానీ $\Th{Q}$ యొక్క 4వ స్వీకృతం ప్రకారం $\eq[(y' + 0)][y']$; 2వ స్వీకృతం ప్రకారం $\eq/[y'][0]$---ఇది వైరుధ్యం. కాబట్టి $\lforall[y][\eq/[(y' +x)][0]]$. $x$, $0$ కాకపోతే 3వ స్వీకృతం వల్ల $\eq[x][z']$ అయ్యే $z$ ఒకటి ఉంది. అప్పుడు $\eq[(y' + z')][0]$ వస్తుంది. 5వ స్వీకృతం వల్ల $\eq[(y' + z)'][0]$ వస్తుంది; అది మళ్లీ 2వ స్వీకృతానికి విరుద్ధం.

2వ భాగానికి $n$పై ఆగమన పద్ధతి వాడండి. $n =0$ అయ్యే ఆధార సందర్భాన్ని చూద్దాం. అప్పుడు $x < \num 1 \lif \eq[x][\num
  0]$ అని చూపాలి. $x < \num 1$ అనుకుందాం. $<$ నిర్వచించే స్వీకృతం వల్ల $\lexists[y][\eq[(y'+x)][0']]$ వస్తుంది. $y$కు ఆ ధర్మం ఉందనుకుందాం; అప్పుడు $\eq[y'+x][0']$.

$\eq[x][0]$ అని చూపాలి. 3వ స్వీకృతం ప్రకారం $x$, 0 కాకపోతే ఏదో $z$కు $z'$తో సమానం. అప్పుడు $\eq[(y' + z')][0']$. $\Th{Q}$ యొక్క 5వ స్వీకృతం ప్రకారం $\eq[(y'+z)'][0']$. 1వ స్వీకృతం ప్రకారం $\eq[(y' +
    z)][0]$. కానీ నిర్వచనం ప్రకారం దీని అర్థం $z < 0$; ఇది 1వ భాగానికి విరుద్ధం.
\end{proof}

\begin{explain}
గణనీయ ప్రమేయాల సమితిని $\Th{Q}$లో ప్రతినిధీకరించగల ప్రమేయాల సమితిగా వర్ణించవచ్చని చూపాం. నిజానికి నిరూపణ మరింత సాధారణమైనది. ప్రతినిధీకరణ నిర్వచనం నుంచి చూస్తే $\Th{Q}$ను విస్తరించే (లేదా $\Th{Q}$ను అన్వయించగల) ఏ సిద్ధాంతమైనా గణనీయ ప్రమేయాలను ప్రతినిధీకరించగలదని గ్రహించడం కష్టం కాదు. మరోవైపు, నిరూపణ భావన గణనీయమైన ఏ \tetoken{వ్యుత్పత్తి}{!!{derivation}} వ్యవస్థలోనైనా ప్రతినిధీకరించగల ప్రతి ప్రమేయం గణనీయమే. అందువల్ల గణనీయ ప్రమేయాల సమితిని పియానో అంకగణితంలో, లేదా జెర్మెలో--ఫ్రెంకెల్ సమితి సిద్ధాంతంలో ప్రతినిధీకరించబడే ప్రమేయాల సమితిగా కూడా వర్ణించవచ్చు. గోడెల్ గుర్తించినట్లు ఇది కొంత ఆశ్చర్యకరం. నిరూపణీయత విషయంలో, ఏ సిద్ధాంతాన్ని తీసుకుంటామనే దానిపై ప్రశ్నలు చాలా సున్నితంగా ఆధారపడతాయని చూస్తాం; సాధారణంగా స్వీకృతాలు ఎంత బలంగా ఉంటే అంత ఎక్కువ నిరూపించవచ్చు. కానీ విస్తృత శ్రేణి స్వీకృత సిద్ధాంతాల్లో ప్రతినిధీకరించగల ప్రమేయాలు సరిగ్గా గణనీయ ప్రమేయాలే.
\end{explain}

\end{document}
